\section{Модель и обучение}

\begin{frame}[c]{Модель и алгоритм обучения}
  \begin{columns}[T]
    \begin{column}{0.47\textwidth}
      \begin{block}{Линейная модель цены}
        Для квартиры предположим: цена зависит от площади линейно. Это упрощение, которое нужно проверить по данным.
      \end{block}
    \end{column}
    \begin{column}{0.47\textwidth}
      \begin{block}{Предсказательная модель}
        \small
        \centerline{$\widehat y=f(x;\theta)$}
        \smallskip
        $x$ --- признаки объекта, $\widehat y$ --- предсказание, $\theta$ --- параметры, которые подбираются по данным.
      \end{block}
    \end{column}
  \end{columns}
\end{frame}

\begin{frame}[c]{Простейшая модель: одна константа}
  \small Всем квартирам предсказываем одну и ту же цену $c$, то есть $f(x)=c$.
  Цены пяти квартир, млн руб.: \textbf{5, 6, 7, 8, 24}.

  \smallskip
  \begin{columns}[T]
    \begin{column}{0.47\textwidth}
      \begin{block}{MAE}
        \centerline{$\displaystyle \mathrm{MAE}=\frac1N\sum_{i=1}^{N}\lvert \widehat y_i-y_i\rvert$}
        \smallskip
        Средний модуль ошибки: насколько в среднем промахиваемся.

        \smallskip
        \lead{Лучшая $c$ --- медиана: 7}
      \end{block}
    \end{column}
    \begin{column}{0.47\textwidth}
      \begin{block}{MSE}
        \centerline{$\displaystyle \mathrm{MSE}=\frac1N\sum_{i=1}^{N}\bigl(\widehat y_i-y_i\bigr)^2$}
        \smallskip
        Средний квадрат ошибки: большие промахи штрафуются сильнее.

        \smallskip
        \lead{Лучшая $c$ --- среднее: 10}
      \end{block}
    \end{column}
  \end{columns}

  \medskip
  {\footnotesize\color{acGray}
  Медиана --- число посередине, если выстроить цены по порядку. Одна дорогая квартира в 24 млн
  сдвинула среднее вверх, а медиану почти не тронула. Здесь $N$ --- число квартир, $i$ --- номер квартиры.\par}
\end{frame}

\begin{frame}{Что нужно, чтобы обучить модель}
  \begin{description}
    \item[Модель] Формула с параметрами $\theta$: меняем параметры --- меняются предсказания.
    \item[Функция потерь] $L(\theta)$ измеряет ошибку модели на обучающих данных.
    \item[Алгоритм обучения] Подбирает $\theta$, уменьшая потери.
  \end{description}
  Например, шаг градиентного спуска:
  \[
    \theta \leftarrow \theta-\eta\,\nabla L(\theta).
  \]
  $\eta$ --- размер шага; $\nabla L$ указывает направление наиболее быстрого роста потерь.
  Параметры меняем в противоположную сторону.
\end{frame}

\begin{frame}[c]{Параметры и гиперпараметры}
  \begin{columns}[T]
    \begin{column}{0.47\textwidth}
      \begin{block}{Параметры}
        Числа внутри модели. Их подбирает алгоритм обучения по данным.
        \begin{itemize}
          \item коэффициенты $k$ и $b$ в прямой $y=kx+b$
          \item веса нейросети
          \item параметры $\theta$ из прошлого слайда
        \end{itemize}
      \end{block}
    \end{column}
    \begin{column}{0.47\textwidth}
      \begin{block}{Гиперпараметры}
        Настройки самого обучения. Их задаёт человек до начала обучения.
        \begin{itemize}
          \item шаг $\eta$ в градиентном спуске
          \item степень многочлена
          \item число слоёв нейросети
          \item сколько раз пройти по данным
        \end{itemize}
      \end{block}
    \end{column}
  \end{columns}
\end{frame}

\begin{frame}{Обучающая, валидационная и тестовая выборки}
  \begin{tabularx}{\linewidth}{@{}l Y@{}}
    \toprule
    \textbf{Выборка} & \textbf{Для чего используется}\\
    \midrule
    Обучающая (train) & Подбираем параметры модели.\\
    Валидационная & Выбираем гиперпараметры и сравниваем модели.\\
    Тестовая (test) & Один раз оцениваем выбранное решение.\\
    \bottomrule
  \end{tabularx}\par
  \begin{block}{Не допускать утечки данных}
    Разделяем данные до обучения. Преобразования, например заполнение пропусков средним, настраиваем только на обучающей части.
  \end{block}
  Для временных данных проверяем будущее по прошлому; связанные объекты не распределяем случайно по разным частям.
\end{frame}

\begin{frame}[c]{Недообучение и переобучение}
  \footnotesize\color{acGray}
  \textcolor{acOrange}{$\bullet$}~обучающие точки \qquad $\circ$~отложенные точки

  \smallskip
  \begin{columns}[T]
    \begin{column}{0.31\textwidth}
      \begin{block}{Слишком простая}
        \centering% underfit: straight line (degree 1).
% Least-squares fit (numpy.polyfit) to the 8 orange training points;
% hollow points are held-out test points. Data: y = 4(x-0.5)^2 + 0.15 + 0.25x
% + N(0, 0.2) noise, numpy seed 7.
% train MSE = 0.0992, test MSE (3 pts) = 0.0906
\begin{tikzpicture}[x=3.38cm,y=1.35cm]
  \draw[chart axis] (0,-0.12) -- (1.06,-0.12);
  \draw[chart axis] (0,-0.12) -- (0,1.42);
  \node[chart label,anchor=north east] at (1.06,-0.12) {$x$};
  \node[chart label,anchor=north east] at (0,1.42) {$y$};
  \begin{scope}
    \clip (0.01,-0.10) rectangle (1.04,1.40);
    \draw[chart curve] plot[smooth] coordinates {(0.025,0.401) (0.035,0.404) (0.044,0.408) (0.054,0.412) (0.063,0.416) (0.073,0.420) (0.083,0.424) (0.092,0.428) (0.102,0.432) (0.111,0.436) (0.121,0.440) (0.131,0.443) (0.140,0.447) (0.150,0.451) (0.159,0.455) (0.169,0.459) (0.179,0.463) (0.188,0.467) (0.198,0.471) (0.207,0.475) (0.217,0.478) (0.227,0.482) (0.236,0.486) (0.246,0.490) (0.255,0.494) (0.265,0.498) (0.275,0.502) (0.284,0.506) (0.294,0.510) (0.303,0.513) (0.313,0.517) (0.323,0.521) (0.332,0.525) (0.342,0.529) (0.351,0.533) (0.361,0.537) (0.371,0.541) (0.380,0.545) (0.390,0.549) (0.399,0.552) (0.409,0.556) (0.419,0.560) (0.428,0.564) (0.438,0.568) (0.447,0.572) (0.457,0.576) (0.467,0.580) (0.476,0.584) (0.486,0.587) (0.495,0.591) (0.505,0.595) (0.515,0.599) (0.524,0.603) (0.534,0.607) (0.543,0.611) (0.553,0.615) (0.563,0.619) (0.572,0.623) (0.582,0.626) (0.591,0.630) (0.601,0.634) (0.611,0.638) (0.620,0.642) (0.630,0.646) (0.639,0.650) (0.649,0.654) (0.659,0.658) (0.668,0.661) (0.678,0.665) (0.687,0.669) (0.697,0.673) (0.707,0.677) (0.716,0.681) (0.726,0.685) (0.735,0.689) (0.745,0.693) (0.755,0.696) (0.764,0.700) (0.774,0.704) (0.783,0.708) (0.793,0.712) (0.803,0.716) (0.812,0.720) (0.822,0.724) (0.831,0.728) (0.841,0.732) (0.851,0.735) (0.860,0.739) (0.870,0.743) (0.879,0.747) (0.889,0.751) (0.899,0.755) (0.908,0.759) (0.918,0.763) (0.927,0.767) (0.937,0.770) (0.947,0.774) (0.956,0.778) (0.966,0.782) (0.975,0.786) (0.985,0.790)};
  \end{scope}
  \node[chart test] at (0.230,0.230) {};
  \node[chart test] at (0.480,0.180) {};
  \node[chart test] at (0.870,0.535) {};
  \node[chart train] at (0.065,0.825) {};
  \node[chart train] at (0.176,0.490) {};
  \node[chart train] at (0.311,0.469) {};
  \node[chart train] at (0.409,0.357) {};
  \node[chart train] at (0.542,0.314) {};
  \node[chart train] at (0.675,0.255) {};
  \node[chart train] at (0.780,0.653) {};
  \node[chart train] at (0.953,1.348) {};
\end{tikzpicture}%
\par
        Прямая не ловит изгиб данных.
      \end{block}
    \end{column}
    \begin{column}{0.31\textwidth}
      \begin{block}{Подходящая}
        \centering% good fit: parabola (degree 2).
% Least-squares fit (numpy.polyfit) to the 8 orange training points;
% hollow points are held-out test points. Data: y = 4(x-0.5)^2 + 0.15 + 0.25x
% + N(0, 0.2) noise, numpy seed 7.
% train MSE = 0.0102, test MSE (3 pts) = 0.0780
\begin{tikzpicture}[x=3.38cm,y=1.35cm]
  \draw[chart axis] (0,-0.12) -- (1.06,-0.12);
  \draw[chart axis] (0,-0.12) -- (0,1.42);
  \node[chart label,anchor=north east] at (1.06,-0.12) {$x$};
  \node[chart label,anchor=north east] at (0,1.42) {$y$};
  \begin{scope}
    \clip (0.01,-0.10) rectangle (1.04,1.40);
    \draw[chart curve] plot[smooth] coordinates {(0.025,0.991) (0.035,0.958) (0.044,0.926) (0.054,0.895) (0.063,0.864) (0.073,0.835) (0.083,0.806) (0.092,0.777) (0.102,0.750) (0.111,0.723) (0.121,0.697) (0.131,0.672) (0.140,0.647) (0.150,0.624) (0.159,0.601) (0.169,0.578) (0.179,0.557) (0.188,0.536) (0.198,0.516) (0.207,0.496) (0.217,0.478) (0.227,0.460) (0.236,0.443) (0.246,0.427) (0.255,0.411) (0.265,0.396) (0.275,0.382) (0.284,0.369) (0.294,0.356) (0.303,0.344) (0.313,0.333) (0.323,0.322) (0.332,0.313) (0.342,0.304) (0.351,0.296) (0.361,0.288) (0.371,0.281) (0.380,0.275) (0.390,0.270) (0.399,0.266) (0.409,0.262) (0.419,0.259) (0.428,0.257) (0.438,0.255) (0.447,0.254) (0.457,0.254) (0.467,0.255) (0.476,0.257) (0.486,0.259) (0.495,0.262) (0.505,0.265) (0.515,0.270) (0.524,0.275) (0.534,0.281) (0.543,0.288) (0.553,0.295) (0.563,0.303) (0.572,0.312) (0.582,0.322) (0.591,0.332) (0.601,0.343) (0.611,0.355) (0.620,0.368) (0.630,0.381) (0.639,0.395) (0.649,0.410) (0.659,0.425) (0.668,0.442) (0.678,0.459) (0.687,0.477) (0.697,0.495) (0.707,0.514) (0.716,0.535) (0.726,0.555) (0.735,0.577) (0.745,0.599) (0.755,0.622) (0.764,0.646) (0.774,0.670) (0.783,0.696) (0.793,0.721) (0.803,0.748) (0.812,0.776) (0.822,0.804) (0.831,0.833) (0.841,0.862) (0.851,0.893) (0.860,0.924) (0.870,0.956) (0.879,0.989) (0.889,1.022) (0.899,1.056) (0.908,1.091) (0.918,1.127) (0.927,1.163) (0.937,1.200) (0.947,1.238) (0.956,1.277) (0.966,1.316) (0.975,1.356) (0.985,1.397)};
  \end{scope}
  \node[chart test] at (0.230,0.230) {};
  \node[chart test] at (0.480,0.180) {};
  \node[chart test] at (0.870,0.535) {};
  \node[chart train] at (0.065,0.825) {};
  \node[chart train] at (0.176,0.490) {};
  \node[chart train] at (0.311,0.469) {};
  \node[chart train] at (0.409,0.357) {};
  \node[chart train] at (0.542,0.314) {};
  \node[chart train] at (0.675,0.255) {};
  \node[chart train] at (0.780,0.653) {};
  \node[chart train] at (0.953,1.348) {};
\end{tikzpicture}%
\par
        Описывает зависимость и на обучающих, и на отложенных точках.
      \end{block}
    \end{column}
    \begin{column}{0.31\textwidth}
      \begin{block}{Переобученная}
        \centering% overfit: polynomial of degree 7.
% Least-squares fit (numpy.polyfit) to the 8 orange training points;
% hollow points are held-out test points. Data: y = 4(x-0.5)^2 + 0.15 + 0.25x
% + N(0, 0.2) noise, numpy seed 7.
% train MSE = 0.0000, test MSE (3 pts) = 0.3830
\begin{tikzpicture}[x=3.38cm,y=1.35cm]
  \draw[chart axis] (0,-0.12) -- (1.06,-0.12);
  \draw[chart axis] (0,-0.12) -- (0,1.42);
  \node[chart label,anchor=north east] at (1.06,-0.12) {$x$};
  \node[chart label,anchor=north east] at (0,1.42) {$y$};
  \begin{scope}
    \clip (0.01,-0.10) rectangle (1.04,1.40);
    \draw[chart curve] plot[smooth] coordinates {(0.025,1.870) (0.035,1.523) (0.044,1.244) (0.054,1.023) (0.063,0.849) (0.073,0.717) (0.083,0.619) (0.092,0.548) (0.102,0.500) (0.111,0.470) (0.121,0.454) (0.131,0.449) (0.140,0.451) (0.150,0.458) (0.159,0.469) (0.169,0.481) (0.179,0.494) (0.188,0.505) (0.198,0.516) (0.207,0.524) (0.217,0.530) (0.227,0.533) (0.236,0.533) (0.246,0.532) (0.255,0.527) (0.265,0.521) (0.275,0.512) (0.284,0.502) (0.294,0.491) (0.303,0.479) (0.313,0.466) (0.323,0.453) (0.332,0.439) (0.342,0.426) (0.351,0.414) (0.361,0.402) (0.371,0.391) (0.380,0.381) (0.390,0.372) (0.399,0.364) (0.409,0.357) (0.419,0.351) (0.428,0.346) (0.438,0.342) (0.447,0.338) (0.457,0.336) (0.467,0.333) (0.476,0.331) (0.486,0.330) (0.495,0.328) (0.505,0.326) (0.515,0.323) (0.524,0.320) (0.534,0.317) (0.543,0.313) (0.553,0.308) (0.563,0.303) (0.572,0.297) (0.582,0.291) (0.591,0.284) (0.601,0.277) (0.611,0.270) (0.620,0.264) (0.630,0.258) (0.639,0.254) (0.649,0.251) (0.659,0.250) (0.668,0.252) (0.678,0.257) (0.687,0.266) (0.697,0.279) (0.707,0.297) (0.716,0.321) (0.726,0.350) (0.735,0.387) (0.745,0.430) (0.755,0.481) (0.764,0.539) (0.774,0.605) (0.783,0.679) (0.793,0.760) (0.803,0.847) (0.812,0.942) (0.822,1.041) (0.831,1.144) (0.841,1.248) (0.851,1.353) (0.860,1.455) (0.870,1.550) (0.879,1.636) (0.889,1.707) (0.899,1.758) (0.908,1.784) (0.918,1.778) (0.927,1.732) (0.937,1.637) (0.947,1.484) (0.956,1.261) (0.966,0.957) (0.975,0.559) (0.985,0.050)};
  \end{scope}
  \node[chart test] at (0.230,0.230) {};
  \node[chart test] at (0.480,0.180) {};
  \node[chart test] at (0.870,0.535) {};
  \node[chart train] at (0.065,0.825) {};
  \node[chart train] at (0.176,0.490) {};
  \node[chart train] at (0.311,0.469) {};
  \node[chart train] at (0.409,0.357) {};
  \node[chart train] at (0.542,0.314) {};
  \node[chart train] at (0.675,0.255) {};
  \node[chart train] at (0.780,0.653) {};
  \node[chart train] at (0.953,1.348) {};
\end{tikzpicture}%
\par
        Проходит почти через каждую обучающую точку, но между ними и на новых данных ошибается.
      \end{block}
    \end{column}
  \end{columns}
\end{frame}

\begin{frame}[c]{Ошибка на обучении и валидации}
  \begin{columns}[T]
    \begin{column}{0.56\textwidth}
      \centering
      % Train/validation MSE vs polynomial degree 0..5 for the same data as fit-*.tex.
% Train MSE on the 8 training points, validation MSE on 400 held-out points
% (numpy.polyfit; seed 7, noise 0.2).
% train: 0.1126, 0.0992, 0.0102, 0.0050, 0.0041, 0.0011
% validation:  0.1082, 0.1029, 0.0356, 0.0380, 0.0393, 0.0477
\begin{tikzpicture}[x=1.06cm,y=1cm]
  \draw[chart grid] (0,0.000) -- (5.25,0.000);
  \node[chart label,anchor=east] at (0,0.000) {0{,}000};
  \draw[chart grid] (0,0.570) -- (5.25,0.570);
  \node[chart label,anchor=east] at (0,0.570) {0{,}025};
  \draw[chart grid] (0,1.140) -- (5.25,1.140);
  \node[chart label,anchor=east] at (0,1.140) {0{,}050};
  \draw[chart grid] (0,1.710) -- (5.25,1.710);
  \node[chart label,anchor=east] at (0,1.710) {0{,}075};
  \draw[chart grid] (0,2.280) -- (5.25,2.280);
  \node[chart label,anchor=east] at (0,2.280) {0{,}100};
  \draw[chart grid] (0,2.850) -- (5.25,2.850);
  \node[chart label,anchor=east] at (0,2.850) {0{,}125};
  \draw[chart axis] (0,0) -- (5.25,0);
  \draw[chart axis] (0,0) -- (0,3.000);
  \foreach \d in {0,...,5} {\draw[chart axis] (\d,0) -- (\d,-0.07); \node[chart label,anchor=north] at (\d,-0.07) {\d};}
  \draw[acAccent!60, line width=0.6pt, dash pattern=on 2pt off 2pt] (2,0) -- (2,0.812);
  \draw[acInk, line width=1.3pt, line join=round] plot coordinates {(0,2.567) (1,2.261) (2,0.233) (3,0.114) (4,0.093) (5,0.024)};
  \draw[acOrange, line width=1.3pt, line join=round] plot coordinates {(0,2.467) (1,2.346) (2,0.812) (3,0.867) (4,0.895) (5,1.087)};
  \fill[acInk] (0,2.567) circle (1.9pt);
  \fill[acInk] (1,2.261) circle (1.9pt);
  \fill[acInk] (2,0.233) circle (1.9pt);
  \fill[acInk] (3,0.114) circle (1.9pt);
  \fill[acInk] (4,0.093) circle (1.9pt);
  \fill[acInk] (5,0.024) circle (1.9pt);
  \fill[acOrange] (0,2.467) circle (1.9pt);
  \fill[acOrange] (1,2.346) circle (1.9pt);
  \fill[acOrange] (2,0.812) circle (1.9pt);
  \fill[acOrange] (3,0.867) circle (1.9pt);
  \fill[acOrange] (4,0.895) circle (1.9pt);
  \fill[acOrange] (5,1.087) circle (1.9pt);
  \draw[acAccent, line width=1.2pt] (2,0.812) circle (5.5pt);
  \node[font=\scriptsize\bfseries, text=acAccent, anchor=east, inner sep=1pt] at (2-0.17,0.812) {оптимум};
  % legend
  \begin{scope}[shift={(3.35,2.800)}]
    \fill[acCard] (-0.1,0.2) rectangle (1.15,-0.56);
    \draw[acInk, line width=1.3pt] (0,0) -- (0.4,0); \fill[acInk] (0.2,0) circle (1.9pt);
    \node[chart label, text=acText, anchor=west] at (0.48,0) {train};
    \draw[acOrange, line width=1.3pt] (0,-0.36) -- (0.4,-0.36); \fill[acOrange] (0.2,-0.36) circle (1.9pt);
    \node[chart label, text=acText, anchor=west] at (0.48,-0.36) {validation};
  \end{scope}
  \node[chart label, anchor=north] at (2.625,-0.42) {Степень многочлена (сложность модели)};
  \node[chart label, anchor=south, rotate=90] at (-0.78,1.425) {MSE};
\end{tikzpicture}%

    \end{column}
    \begin{column}{0.40\textwidth}
      \small
      \textbf{train} (тёмная линия) падает с ростом сложности, иногда до нуля.

      \smallskip
      \textbf{validation} (оранжевая линия) сначала падает, потом снова растёт.

      \smallskip
      Минимум ошибки на валидации определяет выбранную сложность модели.

      \smallskip
      \footnotesize\color{acGray}Сложность выбираем на валидационной выборке. Тест оставляем для итоговой оценки.
      По горизонтали растёт сложность модели (степень многочлена от 0 до 5), по вертикали --- MSE.
    \end{column}
  \end{columns}
\end{frame}

\begin{frame}{После обучения}
  \begin{description}
    \item[Запуск] Проверяем скорость, память и обработку неверных входных данных.
    \item[Мониторинг] Следим за качеством; при необходимости переобучаем модель на новых данных.
    \item[A/B-тест] Сравниваем версии на случайных группах пользователей.
      Заранее выбираем метрику и длительность эксперимента.
    \item[Изменение данных] Data drift --- изменение распределения признаков;
      concept drift --- изменение связи признаков с ответом.
  \end{description}
\end{frame}

\begin{frame}{Вопросы}
  \begin{enumerate}
    \item Почему высокая accuracy ещё не гарантирует полезную модель?
    \item Чем параметр отличается от гиперпараметра?
    \item Почему нельзя выбирать модель по тестовой выборке?
    \item Что проверите, если качество после запуска упало?
  \end{enumerate}
  \medskip
  На практике будем работать с массивами NumPy и таблицами pandas.
\end{frame}
